\section*{Capítulo \ref{ch:b3:particle-physics} --- Física de partículas}

\begin{solution}{exo:b3:particle-physics:1}
(a) uud: $\tfrac23 + \tfrac23 - \tfrac13 = +1$; udd: $\tfrac23 -
\tfrac13 - \tfrac13 = 0$. (b) $\tfrac23 + \tfrac13 = +1$ (o
antidown carrega $+\tfrac13$). (c) $\pi^- =
\bar{\text{u}}\text{d}$; $\Delta^{++} = $ uuu ($3\times\tfrac23 =
+2$ --- o estado cuja estatística ajudou a forçar a invenção da
cor). (d) O campo de glúons esticado armazena energia
$\propto$ comprimento; muito antes de um \omterm{def:b3:particle-physics:zoo}{quark} ficar livre, a energia
armazenada supera $2m_{\text{q}}c^2$ e a corda estala num par
novo: você arranca um \omterm{def:b3:particle-physics:zoo}{méson}, nunca um \omterm{def:b3:particle-physics:zoo}{quark} nu.
\end{solution}

\begin{solution}{exo:b3:particle-physics:2}
(a) \omterm{def:b3:particle-physics:zoo}{Léptons}: e$^-$, $\nu_\mu$, $\mu^-$; \omterm{def:b3:particle-physics:zoo}{bárions}: p, n; \omterm{def:b3:particle-physics:zoo}{méson}:
$\pi^+$. (b) Só os \omterm{def:b3:particle-physics:zoo}{hádrons}: $\pi^+$, p, n. (c) Estáveis sozinhos:
e$^-$, p e os neutrinos (o nêutron livre vive quinze
minutos; o $\mu$ e o $\pi$ morrem em microssegundos ou menos). (d) Os
neutros: $\nu_\mu$ e n.
\end{solution}

\begin{solution}{exo:b3:particle-physics:3}
(a) Permitido --- decaimento $\beta^-$ padrão (o nêutron livre o
faz, $T_{1/2} \approx \qty{10}{min}$). (b) Proibido pela energia:
os produtos pesam mais que o próton livre (dentro de um núcleo, a \omterm{prop:b3:nuclear-physics:binding}{energia
de ligação} pode pagar, e o decaimento $\beta^+$ ligado acontece). (c)
Proibido: \omterm{prop:b3:particle-physics:conservation}{número bariônico} $1 \to 0$ (e \omterm{prop:b3:particle-physics:conservation}{número leptônico}
$0 \to -1$). (d) Proibido pelo \omterm{prop:b3:particle-physics:conservation}{número leptônico} \emph{de família}:
a muonicidade sumiria e a eletronicidade apareceria --- jamais
observado, e caçado precisamente por isso.
\end{solution}

\begin{solution}{exo:b3:particle-physics:4}
(a) $6\times0.938 \approx \qty{5.6}{GeV}$. (b) Energia cinética de
$m_{\text{p}}c^2 \approx \qty{0.94}{GeV}$ por feixe
($\sqrt s = 4m_{\text{p}}c^2$ de frente). (c) Cerca de $3:1$:
a conservação do momento condena a maior parte da energia de alvo fixo a
um movimento inútil dos destroços para a frente --- os feixes frontais
gastam tudo. (d) A máquina de 1955 foi projetada para vencer com folga
mínima o limiar de \qty{5.6}{GeV}: os orçamentos de aceleradores são
escritos em \omterm{def:b3:relativistic-dynamics:four-momentum}{massa invariante}, e o antipróton chegou no prazo.
\end{solution}

\begin{solution}{exo:b3:particle-physics:5}
(a) $c\tau = \qty{660}{m}$. (b) Um alcance dilatado de
${\sim}\qty{15}{km}$ precisa de $\gamma \sim 23$: $E = \gamma
m_\mu c^2 \sim \qty{2.4}{GeV}$ --- exatamente a escala cósmica. (c)
No referencial dele, o múon vive os singelos \qty{2.2}{\micro s}
enquanto a atmosfera passa correndo, contraída a
${\sim}\qty{650}{m}$: mesma sobrevivência, outro livro-caixa. (d) \omterm{prop:b3:particle-physics:conservation}{Número
bariônico}: não há \omterm{def:b3:particle-physics:zoo}{bárion} leve o bastante, e o múon é um
\omterm{def:b3:particle-physics:zoo}{lépton}; o decaimento dele tem de mudar sabor, o que só a força fraca
faz --- daí microssegundos senhoriais em vez de yoctossegundos de
força forte.
\end{solution}

\begin{solution}{exo:b3:particle-physics:6}
(a) $\text{d} \to \text{u} + \text{e}^- + \bar\nu_{\text{e}}$:
carga $-\tfrac13 = \tfrac23 - 1 + 0$; \omterm{prop:b3:particle-physics:conservation}{número bariônico}
$\tfrac13 = \tfrac13$; número de família eletrônica $0 = 1 - 1$. (b) Um
decaimento de dois corpos fixa a energia do elétron; o
\emph{contínuo} observado significava ou que a conservação de energia
falhava (Bohr chegou a considerar) ou que um terceiro corpo invisível
dividia o orçamento --- o neutrino de Pauli, detectado vinte e cinco anos
depois. (c) $\hbar/m_{\text{W}}c = 197/80400 \approx
\qty{2.5e-3}{fm}$ --- um milésimo de um próton. (d) Em baixas
energias, a interação tem de vencer uma distância que o mediador pesado
dela torna absurdamente curta, e por isso as amplitudes são minúsculas;
perto de \qty{100}{GeV} o W é acessível e a força ``fraca'' mostra
a verdadeira força eletrofraca dela.
\end{solution}

\begin{solution}{exo:b3:particle-physics:7}
(a) A força forte conserva a estranheza, e por isso ela só consegue fazer
o $+1$ e o $-1$ juntos: um káon com o híperon dele. (b) Treze
ordens de grandeza lentos demais para a força forte: a força
fraca, única quebradora de sabor, assina o atestado de óbito.
(c) $10^{13}$: um segundo contra trezentos mil anos.
(d) O nascimento copioso dizia ``forte''; a morte relutante dizia ``não
forte'' --- só um novo número conservado, feito aos pares pela força
forte e quebrado pela fraca, concilia os dois: a
contabilidade \emph{foi} a descoberta.
\end{solution}

\begin{solution}{exo:b3:particle-physics:8}
(a) No referencial de repouso do par, o momento total é nulo: um
fóton nunca pode ter momento nulo, e dois de costas um para o outro podem
--- a conservação escreve a geometria do PET. (b) $E = mc^2 =
0.002\times(3\times10^8)^2 = \qty{1.8e14}{J}$: dois gramas de
aniquilação superam um quilograma de urânio fissionado. (c) De
\omterm{def:b3:nuclear-physics:nuclide}{isótopos} $\beta^+$ feitos em cíclotron --- o $^{18}$F, embutido em
glicose: a própria química do corpo põe a fonte de pósitrons nas
células mais famintas. (d) Um combustível que custa um milhão de vezes o
que rende, e sem minas naturais: a antimatéria é uma calibração de
detector e um motor de contador de histórias, não uma fonte de energia.
\end{solution}

\begin{solution}{exo:b3:particle-physics:9}
(a) O potencial de Higgs é um poço duplo (um chapéu mexicano): o
ponto simétrico $\phi = 0$ é um cume, e o vácuo rolou
para a borda --- o estado fundamental do universo quebrou a própria
simetria, exatamente como o ímã que escolhe o norte. (b) A escolha
quebra só parte da simetria: a combinação que sobrevive
é o eletromagnetismo, e por isso o mediador dele --- o fóton --- fica com
massa nula, enquanto o W e o Z, mediadores das direções quebradas,
comem a rigidez correspondente como massa. (c) A oscilação radial
do \omterm{def:b3:phase-transitions:order}{parâmetro de ordem} --- no ímã, a
flutuação de $|m|$ em torno do valor de equilíbrio dele perto de
$T_{\text{c}}$. (d) Nos primeiros $10^{-11}$ segundos do
universo (\cref{ch:b3:astrophysics}): acima da temperatura
eletrofraca, toda partícula era sem massa e as duas forças eram
uma só.
\end{solution}

\begin{solution}{exo:b3:particle-physics:10}
(a) Cada \qty{26.7}{MeV} $= \qty{4.28e-12}{J}$ entrega dois
neutrinos: fluxo $= 2\times1360/4.28\times10^{-12} \approx
\qty{6.4e14}{m^{-2}.s^{-1}}$. (b) Por uma unha do polegar:
$6\times10^{10}$ por segundo. Passagens de uma vida inteira pelo seu
corpo de ${\sim}\qty{0.03}{m^2}$ e \qty{0.3}{m} de espessura, ao longo de
${\sim}\qty{2.5e9}{s}$: cerca de $5\times10^{22}$, vezes
$10^{-22}\,\unit{m^{-1}}\times\qty{0.3}{m}$: da ordem de um ---
um único neutrino solar tocará você na vida inteira. (c)
Os $\nu_{\text{e}}$ que faltavam tinham oscilado em $\nu_\mu$ e
$\nu_\tau$ no caminho; a oscilação exige diferenças de massa, e portanto
os neutrinos têm massa --- a primeira rachadura de laboratório no
Modelo Padrão de neutrinos sem massa. (d) A montanha filtra
todo o \emph{resto}: só os neutrinos passam, e assim quilômetros de
rocha transformam o céu mais barulhento no laboratório mais silencioso.
\end{solution}

\begin{solution}{exo:b3:particle-physics:11}
(a) $\gamma = 6.8\times10^{12}/9.38\times10^{8} \approx 7250$:
$1 - v/c \approx 1/2\gamma^2 \approx 10^{-8}$ --- três metros
por segundo aquém da luz. (b) $\lambda \approx hc/E \approx
\qty{3e-5}{fm}$: trinta milésimos do raio de um próton --- território
de \omterm{def:b3:particle-physics:zoo}{quarks}. (c) $13.6\,\text{TeV}/938\,\text{MeV} \approx
14500$ massas de próton --- se tudo isso fosse conversível. (d) A
energia fixa, $\gamma_{\text{e}}/\gamma_{\text{p}} =
m_{\text{p}}/m_{\text{e}} \approx 1836$, e assim as perdas do elétron são
$1836^4 \sim 10^{13}$ piores: prótons para os anéis de força bruta,
elétrons (limpos, pontuais) para a precisão em energia mais baixa.
\end{solution}

\begin{solution}{exo:b3:particle-physics:12}
(a) $3\times10^{20}\times1.6\times10^{-19} \approx \qty{48}{J}$:
uma bola de tênis sacada com força --- num só próton. (b) $\gamma \approx
3\times10^{11}$; os $10^{5}$ anos-luz da Galáxia se contraem a
${\sim}10$ \emph{segundos}-luz: ele atravessa a Via Láctea, pelo
relógio dele, em segundos. (c) No referencial do próton, as
micro-ondas mansas da CMB chegam desviadas para o azul, virando raios
$\gamma$ energéticos o bastante para arrancar píons dele: acima do
limiar, o próprio universo é um absorvedor, e prótons assim têm de ser
locais. (d) Uma partícula dessas por quilômetro quadrado por século, de
espécie desconhecida, apontada por ninguém: as $10^{9}$ colisões por
segundo do LHC, de feixes conhecidos a energia conhecida, compram uma
estatística e um controle que presente cósmico algum iguala.
\end{solution}

\begin{solution}{pb:b3:particle-physics:1}
\textbf{1.} Um próton primário (muitas vezes de \unit{GeV}--\unit{TeV})
atinge um núcleo lá em cima (força forte), espalhando píons; os
píons carregados decaem pela força fraca em múons e neutrinos;
os múons, que ionizam só de leve (eletromagnético), correm para o
solo.
\textbf{2.} $c\tau_\pi \approx \qty{8}{m}$: mesmo um $\gamma$ grande
compra aos píons só centenas de metros, e eles também morrem por
\omterm{rem:b3:particle-physics:forces}{interações fortes} no ar; o múon, que não sente a força forte e
vive $85\times$ mais, é quem viaja.
\textbf{3.} $400\,\unit{cm^2}\times1\,\text{min}^{-1} = 400$
por minuto --- cerca de 7 por segundo.
\textbf{4.} Os clarões de ruído e a radioatividade local disparam uma
pá ao acaso; disparar \emph{as duas} dentro de \qty{100}{ns}
exige uma única partícula atravessando o par --- a geometria como
filtro.
\textbf{5.} $R_{\text{acc}} = R_1R_2\Delta t = 100\times100
\times10^{-7} = \qty{e-3}{s^{-1}}$ --- uma contagem falsa por
quarto de hora contra cem múons por minuto: desprezível.
\textbf{6.} Só os múons dentro do cone que enfia as duas
pás (quase verticais, dentro de ${\sim}30$--$40^\circ$) são
aceitos: o par mede uma direção, e não só uma taxa ---
é isso que faz dele um \emph{telescópio}.
\textbf{7.} Os múons inclinados atravessam mais atmosfera: mais perda de
energia e mais \omterm{prop:b3:relativistic-kinematics:dilation}{tempo próprio} em voo, e por isso menos sobrevivem ---
mais ou menos a lei $\cos^2\theta$ que as suas medidas inclinadas traçam.
\textbf{8.} $c\tau = \qty{660}{m}$: um múon clássico vindo de
\qty{15}{km} sobrevive com probabilidade
$\eu^{-15000/660} \approx 10^{-10}$ --- o sótão contaria
um múon por década.
\textbf{9.} $\gamma = 4000/106 \approx 38$.
\textbf{10.} Tempo de vida dilatado $\gamma\tau \approx
\qty{84}{\micro s}$: alcance de ${\sim}\qty{25}{km}$ --- a taxa do
sótão é exatamente o que a relatividade prevê.
\textbf{11.} No referencial do múon, o tempo de vida é os singelos
\qty{2.2}{\micro s}, mas a atmosfera está contraída a
$15000/38 \approx \qty{400}{m}$: atravessada com tempo de sobra.
\textbf{12.} $t_{\text{propr}} \approx 15000/(38\times3\times
10^8) = \qty{1.3}{\micro s}$: sobrevivência $\eu^{-1.3/2.2} \approx
0.55$ --- metade chega, e o seu contador é o censo deles.
\textbf{13.} Os experimentos clássicos de múons na montanha contra o
nível do mar (Rossi--Hall e, em filme, Frisch--Smith): o decaimento
clássico prevê essencialmente contagem nenhuma; as suas centenas por
minuto são a relatividade restrita, verificada acima do isolamento.
\textbf{14.} Um \omterm{def:b3:particle-physics:zoo}{lépton} carregado da segunda geração:
carga $-e$, spin $\tfrac12$, massa de \qty{106}{MeV} --- o
gêmeo pesado do elétron.
\textbf{15.} Carga: $-1 = -1 + 0 + 0$; família eletrônica:
$0 = 1 + (-1) + 0$; família muônica: $1 = 0 + 0 + 1$. Todos os livros
fecham.
\textbf{16.} Três corpos dividem o orçamento em todas as proporções,
e por isso o espectro do elétron é contínuo; ele termina quando os dois
neutrinos recuam juntos contra o elétron: cerca de
$m_\mu c^2/2 \approx \qty{53}{MeV}$.
\textbf{17.} Isso quebraria o \omterm{prop:b3:particle-physics:conservation}{número leptônico} de família --- o único
livro que o Modelo Padrão fecha sem saber por quê. Achá-lo
seria o primo, no setor carregado, das oscilações de neutrinos:
física nova garantida, e é por isso que os experimentos o perseguem até
uma parte em $10^{13}$.
\textbf{18.} A entrada do múon faz um clarão; microssegundos
depois, o elétron do decaimento dele faz um segundo na mesma pá. O
histograma dos atrasos é uma exponencial cuja inclinação \emph{é}
$\tau = \qty{2.2}{\micro s}$: um tempo de vida medido numa prateleira.
\textbf{19.} Todas as fórmulas de Bohr escalam com a massa: a órbita do
hidrogênio muônico encolhe por $207$, até ${\sim}\qty{250}{fm}$ ---
tão fundo dentro da nuvem eletrônica que o múon sente o
tamanho do próton: os átomos muônicos são réguas de prótons.
\textbf{20.} $10^{15}/10^{6} = \qty{e9}{eV}$: cerca de um GeV por
partícula --- os múons do seu sótão são cidadãos médios do chuveiro.
\textbf{21.} Os múons relativísticos correm mais que a luz na água e
brilham no azul de Cherenkov para as fotomultiplicadoras do tanque. A
pegada de um chuveiro de $10^{20}\,\unit{eV}$ cobre muitos quilômetros
quadrados, e assim tanques esparsos a amostram como baldes numa
chuva.
\textbf{22.} Uma partícula deposita energia num meio que a
converte em luz, uma fotomultiplicadora converte a luz em
elétrons, e a eletrônica os conta: todo calorímetro, do
sótão ao colisor, é esta frase.
\textbf{23.} O raio do enrolamento dá o momento, e o sentido dele dá
o sinal da carga; o fóton, o nêutron e o neutrino ---
neutros --- não deixam traço, só as dívidas deles na
contabilidade.
\textbf{24.} A natureza acelera prótons isolados dez milhões de
vezes além do LHC --- humilhante --- mas a um por quilômetro
quadrado por século: as nossas máquinas trocam energia de pico por
$10^{9}$ colisões controladas por segundo.
\textbf{25.} Nascido de um próton dez quilômetros acima; entregue
vivo pela \omterm{prop:b3:relativistic-kinematics:dilation}{dilatação do tempo} ($\gamma \approx 38$, com metade sobrevivendo);
identificado por dois clarões a cem nanossegundos de distância; e
explicado, por inteiro, por uma tabela de doze cadeiras e quatro
forças.
\end{solution}
